\documentclass[11pt,letterpaper]{article}
\usepackage{fontspec}
\usepackage{tocloft}
\usepackage{multicol}
\usepackage{nicefrac}
\usepackage[left=1.3in,right=1.3in,top=1in,bottom=1in]{geometry}
\setmainfont[Scale=1.1,AutoFakeBold=1.5,AutoFakeSlant=0.3]{Doves Type}

% Prevent page breaks within list items
\widowpenalty=10000
\clubpenalty=10000
\interlinepenalty=500

\title{Sheet Pan Salmon with Dijon-Tarragon Beurre Blanc}
\author{}
\date{}

\begin{document}

\maketitle
\thispagestyle{empty}

\begin{quote}\textit{Brussels sprouts and whole garlic cloves roast first for a deep caramelized head start, then asparagus and salmon fillets join the pan in stages so each finishes exactly when it should. A Dijon-tarragon beurre blanc---built on a shallot-wine-lemon reduction, enriched with mashed roasted garlic, and mounted with cold butter---is spooned over everything at the table, with the remaining whole roasted garlic cloves served alongside. Serves 4.}\end{quote}

\section*{Ingredients}
\setlength{\columnsep}{20pt}
\begin{multicols}{2}
\noindent
    Brussels sprouts, halved \dotfill 1~lb. \\
    Garlic cloves, peeled \dotfill 8 \\
    Neutral oil, divided \dotfill 3~Tbsp. \\
    Dried thyme \dotfill 1~tsp. \\
    Kosher salt, divided \dotfill 1\nicefrac{1}{4}~tsp. \\
    Black pepper, divided \dotfill 1~tsp. \\
    Asparagus, trimmed \dotfill 1~lb. \\
    Salmon fillets \dotfill 4 \\
    \columnbreak
    Smoked paprika \dotfill \nicefrac{1}{2}~tsp. \\
    Red pepper flakes \dotfill \nicefrac{1}{4}~tsp. \\
    Shallot, minced \dotfill 3~Tbsp. \\
    Dry white wine \dotfill \nicefrac{3}{4}~cup \\
    Lemon juice \dotfill 3~Tbsp. (from 1~lemon) \\
    Dried tarragon \dotfill 1~tsp. \\
    Dried parsley \dotfill 1~tsp. \\
    Dijon mustard \dotfill 1\nicefrac{1}{2}~Tbsp. \\
    Salted butter, cold, cubed \dotfill 10~Tbsp. \\
\end{multicols}

\section*{Directions}

\noindent
Thaw \textbf{salmon} fully if frozen (overnight in the refrigerator, or sealed in a bag in cold water for \textit{30--45~minutes}) ---
Halve \textbf{Brussels sprouts} through the core; peel \textbf{garlic cloves}; toss both with 2~Tbsp \textbf{neutral oil}, \textbf{dried thyme}, \nicefrac{1}{2}~tsp. \textbf{kosher salt}, and \nicefrac{1}{4}~tsp. \textbf{black pepper} in \textit{Medium Bowl~\#1} ---
Toss \textbf{asparagus} with remaining 1~Tbsp \textbf{neutral oil}, \nicefrac{1}{4}~tsp. \textbf{kosher salt}, and \nicefrac{1}{4}~tsp. \textbf{black pepper} in \textit{Medium Bowl~\#2} ---
Combine \textbf{smoked paprika}, \textbf{red pepper flakes}, \nicefrac{1}{4}~tsp. \textbf{kosher salt}, and \nicefrac{1}{4}~tsp. \textbf{black pepper} in \textit{Small Bowl~\#1} (salmon rub) ---
Mince \textbf{shallot} ---
Measure \textbf{dried tarragon} and \textbf{dried parsley} together in \textit{Small Bowl~\#2} ---
Cube \textbf{salted butter} and keep cold ---
Measure \textbf{wine}, \textbf{Dijon mustard}, and \textbf{lemon juice}; have a little warm water ready

\begin{enumerate}
    \item Preheat oven to \textit{425°F}. Line a rimmed half-sheet pan with parchment or foil, if desired, for easy cleanup.

    \item Arrange \textbf{Brussels sprouts} and \textbf{garlic cloves} (\textit{Medium Bowl~\#1}) in a single layer across the pan, sprouts cut-side down, spacing everything so it isn't touching---crowding traps steam and prevents browning. Roast for \textit{16--18~minutes}, until the sprouts' cut sides are deep golden at the edges and the garlic cloves are soft and lightly browned when pressed. Do not stir.

    \item While vegetables roast, pat \textbf{salmon} fillets dry with paper towels and rub both sides with the spice mixture (\textit{Small Bowl~\#1}).

    \item Remove the pan from the oven. Transfer \textbf{garlic cloves} to \textit{Small Bowl~\#3} and tent loosely with foil to keep warm---pull them now rather than leaving them for the full roast: peeled cloves have no skin left to protect them, and direct high heat past about \textit{20~minutes} turns them bitter instead of sweet. Push \textbf{Brussels sprouts} to the ends of the pan and add \textbf{asparagus} (\textit{Medium Bowl~\#2}) to the open space in the center. Return the pan to the oven. Immediately combine \textbf{shallot}, \textbf{wine}, \textbf{lemon juice}, and the dried tarragon and parsley (\textit{Small Bowl~\#2}) in a small saucepan and bring to a simmer over medium heat. Roast the \textbf{asparagus} for \textit{5--7~minutes}, until it turns bright green and just starts to soften; keep the sauce at a steady simmer.

    \item Add the \textbf{salmon} and finish the reduction:
    \begin{enumerate}
        \item Remove the pan from the oven and add \textbf{salmon} fillets to the center, nestled among the \textbf{asparagus}. Return the pan to the oven and roast for \textit{8--10~minutes}; check \textbf{salmon} starting at \textit{7~minutes}, since fillet thickness varies. Salmon is done when opaque throughout and it reaches \textit{125--130°F} on an instant-read thermometer at the thickest point. Asparagus should be bright green and tender-crisp when pierced with a fork; Brussels sprouts should be deeply caramelized and tender at the center when pierced with a knife. Transfer \textbf{asparagus} to a platter and tent with foil once it is tender-crisp, rather than leaving it in to overcook.
        \item Meanwhile, continue simmering the sauce until reduced to about 4--5~Tbsp. of slightly syrupy liquid, \textit{10--12~minutes} total from when it first reached a simmer: the shallots look glazed, the liquid is concentrated, and the raw alcohol smell is gone. Watch closely near the end, as the last of the liquid evaporates quickly. Add a splash of \textbf{wine} to return a nearly dry pan to 4--5~Tbsp. Remove from heat.
    \end{enumerate}

    \item Transfer \textbf{salmon} to a platter and tent loosely with foil; rest for \textit{3--5~minutes}---internal temperature will rise about \textit{5°F} further. Arrange \textbf{Brussels sprouts} and \textbf{asparagus} alongside. From \textit{Small Bowl~\#3}, reserve 2--3 \textbf{garlic cloves} and mash with a fork; set the rest aside whole to serve.

    \item While \textbf{salmon} rests, finish the sauce:
    \begin{enumerate}
        \item Whisk the mashed \textbf{garlic} and \textbf{Dijon mustard} into the reduction; set the saucepan over very low heat.
        \item Add \textbf{butter} a few cubes at a time, whisking constantly and letting each addition melt and emulsify fully before adding more. Keep the sauce between \textit{122--131°F} throughout---if it looks like it's approaching a simmer, pull the pan off the heat briefly while continuing to whisk. Above \textit{149°F} the emulsion will break, turning the sauce greasy and separated instead of glossy.
        \item Once all the \textbf{butter} is incorporated and the sauce is smooth and glossy, whisk in warm water, 1~tsp. at a time, as needed until the sauce coats the back of a spoon in an even sheet and a line drawn through it with a finger holds clearly (napp\'{e}), and the wine and tarragon aroma bloom---stop by about 1~Tbsp. Taste and adjust with the remaining \textbf{kosher salt} and \textbf{black pepper}.
    \end{enumerate}

    \item Spoon the sauce generously over \textbf{salmon} and vegetables. Serve immediately with the reserved whole roasted \textbf{garlic cloves} alongside.
\end{enumerate}

\newpage

{\footnotesize
\setlength{\columnsep}{20pt}
\setlength{\multicolsep}{6pt}
\begin{multicols}{2}
\setlength{\parindent}{0pt}
\setlength{\parskip}{2pt}
\setlength{\itemsep}{0pt}
\setlength{\parsep}{0pt}

\subsection*{Yield}
\begin{itemize}
    \item Serves 4 (4 fillets, 3--4~oz. each, plus 1~lb. asparagus, 1~lb. Brussels sprouts, and about \nicefrac{3}{4}--1~cup sauce)
\end{itemize}

\subsection*{Equipment Required}
\begin{itemize}
    \item Large rimmed half-sheet pan (18$\times$13-inch)
    \item Parchment paper or foil (optional, for lining)
    \item Instant-read thermometer
    \item Small saucepan (for the sauce)
    \item Whisk
    \item Fork (for mashing garlic and checking vegetable tenderness)
    \item Small prep bowls (3), Medium prep bowls (2)
    \item Paper towels
    \item Measuring cups and spoons
    \item Citrus juicer (optional)
\end{itemize}

\subsection*{Mise en Place}
\begin{itemize}
    \item Medium Bowl \#1 --- \textbf{Brussels sprouts} and \textbf{garlic cloves} tossed with \textbf{neutral oil}, \textbf{dried thyme}, \textbf{salt}, \textbf{pepper}
    \item Medium Bowl \#2 --- \textbf{asparagus} tossed with \textbf{neutral oil}, \textbf{salt}, \textbf{pepper}
    \item Small Bowl \#1 --- salmon rub: \textbf{smoked paprika}, \textbf{red pepper flakes}, \textbf{salt}, \textbf{pepper}
    \item Small Bowl \#2 --- \textbf{dried tarragon} and \textbf{dried parsley}, measured together for the sauce
    \item Small Bowl \#3 --- reserved roasted \textbf{garlic cloves} (filled partway through cooking, not during prep)
    \item Mince \textbf{shallot}; cube \textbf{butter} and keep cold; measure \textbf{wine}, \textbf{Dijon mustard}, and \textbf{lemon juice}; have a little warm water ready
    \item Thaw \textbf{salmon} fully before cooking; pat dry before seasoning
\end{itemize}

\subsection*{Ingredient Tips}
\begin{itemize}
    \item \textbf{Salmon}: Frozen, skinless, boneless, 3--4~oz. fillets; thaw fully and pat dry before seasoning so they roast rather than steam
    \item \textbf{Brussels sprouts}: Choose similarly sized sprouts (1--1\nicefrac{1}{2}~inch) for even cooking
    \item \textbf{Asparagus}: Medium-thick, pencil-width spears roast most evenly; thin spears may finish early
    \item \textbf{Smoked paprika}: Use sweet/mild, not hot---\textbf{red pepper flakes} already supply the heat
    \item \textbf{Dried tarragon \& parsley}: Added to the reduction to hydrate; substitute fresh herbs at 3$\times$ the amount, stirred in raw off heat at the end
    \item \textbf{Salted butter}: Sodium varies by brand---no salt is added to the sauce base; taste after the butter is incorporated and the sauce is loosened to napp\'{e}
    \item \textbf{Dry white wine}: Sauvignon Blanc, Pinot Grigio, or similar---use one you would drink; avoid heavily oaked Chardonnay
    \item \textbf{Shallot}: Yellow onion substitutes in a pinch, but shallot's milder sweetness is preferred
\end{itemize}

\subsection*{Preparation Tips}
\begin{itemize}
    \item Pull \textbf{garlic} off the pan once softened---peeled cloves scorch past about \textit{20~minutes} of direct heat
    \item Don't crowd the pan at any stage---steam trapped between ingredients prevents browning
    \item Start the sauce as soon as the \textbf{asparagus} goes in so the reduction can reach 4--5~Tbsp. while the \textbf{salmon} roasts
    \item Keep \textbf{butter} cold and add it gradually to control the sauce's temperature; if it breaks, whisking in one more cold cube off heat can help
    \item A splash of warm water at the end loosens a tight sauce to napp\'{e} and opens the flavor---fat holds the wine and tarragon aromatics until water gives them a way out. Use water, not more wine. Ice-cold water is for a broken sauce, not this finish
    \item Pull \textbf{salmon} at \textit{125--130°F} for moist texture; it rises about \textit{5°F} further while resting. USDA recommends \textit{145°F}; commercially frozen salmon is generally considered safe at the lower temperature used here
    \item Make the sauce fresh, close to serving---it doesn't reheat well; hold briefly off direct heat only
\end{itemize}

\subsection*{Make Ahead \& Storage}
\begin{itemize}
    \item \textbf{Brussels sprouts} and \textbf{asparagus} can be trimmed up to 2~days ahead and refrigerated
    \item The salmon rub (\textit{Small Bowl~\#1}) can be mixed several days ahead
    \item Make the sauce fresh---it breaks if reheated
    \item Leftover salmon and vegetables keep \textit{2~days} refrigerated; reheat gently at \textit{300°F}
    \item Leftover roasted \textbf{garlic} keeps \textit{3--4~days} refrigerated; good mashed onto bread
\end{itemize}

\subsection*{Serving Suggestions}
\begin{itemize}
    \item Serve with rice cooked in chicken stock, bay leaf, and a pinch of \textbf{dried thyme} (echoing the sprouts), alongside prepared mashed potatoes
    \item Spoon extra sauce over the rice or mashed potatoes as well as the salmon
    \item Set the reserved whole roasted \textbf{garlic cloves} on the table for guests to mash onto their own plates
\end{itemize}

\end{multicols}
}

\end{document}
